% UNUSED since 2026-10-22: the method paper folds this section into sections/method.tex (Sec. 3.1).
\section{Interaction in Physical Time}
\label{sec:formulation}

\begin{figure*}[t]
\centering
\includegraphics[width=0.76\textwidth,trim=0 0 0 26,clip]{interaction_timing.pdf}
\caption{\textbf{Three consequences of acting in physical time} (schematic). (a)
Two histories reach the same pose at capture time $\tau_t$ but imply different
motions. (b) At the same application time $\tapp$, they lead to different target
poses, in translation and in rotation. (c) A contact requires the end-effector
frame $E$ to be aligned with the target $O$ within $\varepsilon$, whereas the
context matters only through the gap it leaves, with a margin $m\gg\varepsilon$.}
\label{fig:principles}
\end{figure*}

\paragraph{Physical time.}
We consider a policy that receives images and robot states together with their
capture times. At
query $t$, its history is $\mathcal H_t=\{(o_i,s_i,\tau_i)\}_{\tau_i\le\tau_t}$,
and it outputs a chunk $A_t=(\vect a_{t,0},\ldots,\vect a_{t,H-1})$ whose $j$-th
command takes effect at
\begin{equation}
\tapp_{t,j}=\tau_t+\ell_t+j\Delta_a ,
\label{eq:apptime}
\end{equation}
where $\ell_t$ is the latency from capture to the first application
(preprocessing, inference, queuing, and actuation) and $\Delta_a$ is the control
period. We call $\Delta_{t,j}=\tapp_{t,j}-\tau_t$ the \emph{lead} of command $j$,
which must fit the scene at $\tapp_{t,j}$ rather than at $\tau_t$.

\paragraph{Target and context.}
Let $W_t(\vect x,\tau)\in\R^3$ be the displacement between $\tau_t$ and $\tau$ of
the scene point that is at $\vect x$ at time $\tau_t$, expressed in the robot base
frame. We split the scene by the precision that the task requires of it
(Fig.~\ref{fig:principles}). The \emph{target} $\mathcal O$ is the rigid body
whose pose error enters a contact or containment tolerance $\varepsilon$, e.g., an
object to grasp or a receptacle to place into. The \emph{context} $\mathcal C$
contains the bodies that only need to be avoided---a belt, a wire, an
occluder---with clearance $m$.
On the target, the motion field is set by its pose
$T(\tau)=(R(\tau),\vect p(\tau))\in\mathrm{SE}(3)$:
\begin{equation}
W_t(\vect x,\tau)=T(\tau)\,T(\tau_t)^{-1}\vect x-\vect x,\qquad \vect x\in\mathcal O .
\label{eq:rigidfield}
\end{equation}

\paragraph{An interaction error budget.}
Consider a contact point at lever arm $r$ from the target origin and a policy that
acts, explicitly or implicitly, on an estimate $\hat T$ of the target pose at
$\tapp$. With a rotation error $\theta$, the contact error satisfies
\begin{equation}
e_{\rm c}\le\|\hat{\vect p}(\tapp)-\vect p(\tapp)\|+2r\sin(\theta/2).
\label{eq:contact}
\end{equation}
We write the true motion over the lead $\Delta=\tapp-\tau_t$ as
$\vect p(\tapp)=\vect p(\tau_t)+\vect v\Delta+\boldsymbol\rho$, where
$\boldsymbol\rho$ is what constant-velocity extrapolation misses (acceleration,
rolling, impacts), and the estimate as
$\hat{\vect p}(\tapp)=\hat{\vect p}(\tau_t)+\hat{\vect v}\hat\Delta+\hat{\boldsymbol\rho}$.
Adding and subtracting $\vect v\hat\Delta$ gives
\begin{equation}
\|\hat{\vect p}-\vect p\|\le
\underbrace{e_{\rm loc}}_{\text{now}}
+\underbrace{\hat\Delta\,e_{v}}_{\text{lead}\times\text{motion}}
+\underbrace{\|\vect v\|\,e_{\tau}}_{\text{speed}\times\text{timing}}
+\underbrace{e_{\rm dyn}}_{\text{dynamics}},
\label{eq:budget}
\end{equation}
with $e_{\rm loc}=\|\hat{\vect p}(\tau_t)-\vect p(\tau_t)\|$,
$e_v=\|\hat{\vect v}-\vect v\|$, $e_\tau=|\hat\Delta-\Delta|$, and
$e_{\rm dyn}=\|\hat{\boldsymbol\rho}-\boldsymbol\rho\|$. Rotation decomposes in the
same way (Appendix~\appref{app:analysis}). The context enters differently. If
the commanded path keeps a clearance $m$ from the predicted context positions
$\hat{\vect y}_t(\vect x,\tau)$, a collision requires a prediction error of at
least $m$ at some point and time in the set $\mathcal P$ that the path passes:
\begin{equation}
\text{collision}\ \Rightarrow\ \max_{(\vect x,\tau)\in\mathcal P}
\big\|\hat{\vect y}_t(\vect x,\tau)-\vect y_t(\vect x,\tau)\big\|\ge m .
\label{eq:context}
\end{equation}
A trial succeeds if $e_{\rm c}\le\varepsilon$ and no collision occurs.
Eqs.~\eqref{eq:contact}--\eqref{eq:context} follow from an identity and the
triangle inequality. We use them as a budget rather than as a model of success.

\paragraph{Predictions.}
Given each task's tolerance and speeds (Appendix~\appref{app:roles}): (i) a
policy that ignores motion incurs an error of $\|\vect v\|\Delta$, so it should
fail beyond the critical lead $\ell^\star=\varepsilon/\|\vect v\|$, and its success
should depend on the task, speed, and latency only through
$\|\vect v\|\Delta/\varepsilon$. (ii) A
least-squares fit to $n$ poses with noise $\sigma$ over a span $T$ estimates the
velocity with an error of about $\sigma\sqrt{12/n}/T$ \cite{barshalom2001}, so a
pose history remains useful up to a noise level of
$\sigma^\star=\varepsilon T\sqrt n/(\Delta\sqrt{12})$. (iii) A forecast that is
misaligned with the application times by $\delta$ adds $\|\vect v\|\delta$, which
grows with speed and vanishes for stationary targets. (iv) A context forecast with
path error $e_F$ is sufficient as long as $m>e_F$.

% Protocol table: quantitative predictions of the interaction error budget (Sec. 3), tested in the section Testing the Predictions. UNUSED since 2026-10-22 (method paper; scaling relations in App. G).
\begin{table}[t]
\centering\footnotesize
\caption{\tablelang{\textbf{Predictions of the error budget.} Each row concerns a term of
Eqs.~\eqref{eq:budget}--\eqref{eq:context}, and Sec.~\ref{sec:principles} tests it on held-out \bench scenes by
sweeping the named factor. $\ell^\star$ and $\sigma^\star$ are the critical lead and noise level, and $e_F$ is
the path error of the context flow.}{\textbf{误差预算给出的预测。}每行对应式~\eqref{eq:budget}--\eqref{eq:context} 中的一项，
第~\ref{sec:principles}~节在 \bench 留出场景上扫描其中提到的因素来检验。$\ell^\star$ 与 $\sigma^\star$ 为临界领先时长与临界噪声，$e_F$ 为环境流在路径上的误差。}}
\label{tab:principles}
\setlength{\tabcolsep}{2.8pt}
\renewcommand{\arraystretch}{1.1}
\begin{tabularx}{\columnwidth}{@{}l>{\raggedright\arraybackslash}p{0.86cm}>{\raggedright\arraybackslash}X@{}}
\toprule
\tablelang{Principle}{原则} & \tablelang{Term}{误差项} & \tablelang{Prediction}{预测} \\
\midrule
\tablelang{All}{整体} & $e_{\rm c}/\varepsilon$ & \tablelang{Without compensation, success depends on task, speed, and latency only
through $\|\vect v\|\Delta/\varepsilon$ and drops once it exceeds about 1.}{无补偿时，成功率仅经由 $\|\vect v\|\Delta/\varepsilon$ 随任务、速度和延迟变化，并在该比值超过约 1 后下降。} \\
\tablelang{Remember}{记住} & $\hat\Delta e_v$ & \tablelang{SE(3) history helps beyond $\ell^\star$ (\fillin{X}--\fillin{X}\,ms across
tasks), increasingly with latency and speed.}{SE(3) 历史在超过 $\ell^\star$（各任务 \fillin{X}--\fillin{X}\,ms）后起作用，且随延迟与速度增大。} \\
 & $e_v(\sigma)$ & \tablelang{Its gain persists up to pose noise $\sigma^\star$ (\fillin{X}--\fillin{X}\,mm).}{其增益在位姿噪声达到 $\sigma^\star$（\fillin{X}--\fillin{X}\,mm）前保持。} \\
\tablelang{Predict}{预测} & $e_{\rm dyn}$, $r\theta$ & \tablelang{Learned SE(3) futures beat extrapolation on non-constant
motion and point tracks under rotation.}{学习的 SE(3) 未来在非匀速运动上优于外推，在旋转目标上优于点轨迹。} \\
 & $m$ & \tablelang{Coarse flow suffices down to clearance $e_F$.}{间隙不低于 $e_F$ 时，粗粒度流已足够。} \\
\tablelang{Act}{行动} & $\|\vect v\|e_\tau$ & \tablelang{A 50\,ms horizon shift costs 5\,mm at 0.10\,m/s and nothing for static
targets. Leads in control steps fail when the latency changes.}{在 0.10\,m/s 时 50\,ms 的时域偏移带来 5\,mm 误差，静止目标不受影响。按控制步计的领先时长在延迟变化时失效。} \\
\bottomrule
\end{tabularx}
\end{table}


\paragraph{Design implications.}
Table~\ref{tab:principles} maps each term to the part of a policy that can reduce
it. \emph{Remember.} Given a reliable current pose, a history reduces $e_v$ and,
through acceleration, part of $e_{\rm dyn}$. Because $e_v$ is multiplied by the
lead, the history matters through the target's metric motion state, and pixels,
flow, or video tokens help only if a policy recovers this state from them.
\emph{Predict.} Reducing $e_{\rm dyn}$ requires a model of how the target will
move, and rotation errors are amplified by $r$. The target's future should
therefore be a full SE(3) trajectory, which by Eq.~\eqref{eq:rigidfield}
determines the motion of all its points. Context errors are tolerated up to $m$,
so a 3D flow with an error below $m$ suffices, and appearance does not enter the
budget. \emph{Act.} Every term is evaluated at $\tapp$, which depends on the
latency. Forecasts and commands should therefore be indexed by physical time, with
each command's lead known to the policy. Whether the action should read a forecast
online or be trained to encode it is left open by the budget, and we answer this
question experimentally (Sec.~\ref{sec:principles}).
